\begin{abstract}
Рассматриваются многоканальные системы массового обслуживания с пуассоновскими поступлениями и отрицательными воздействиями. Воздействие выбирает один занятый канал либо одновременно затрагивает всю систему. Исследуются удаление заявок и повторное обслуживание с новой или с первоначально выбранной длительностью. Для фазового обслуживания описано совместное состояние каналов и построены уравнения моментов ожидания и пребывания помеченной заявки. Время пребывания при выборочных воздействиях рассчитывается с учётом изменения числа занятых каналов. Для сохранённой работы используется конечная аппроксимация процесса с непрерывными координатами; зависимость каналов при общих прерываниях сохраняется. Отдельно проверяются фазовая аппроксимация, усечение пространства состояний, экстраполяция и статистическая неопределённость имитации. Приведены результаты для нескольких нагрузок, интенсивностей воздействий и числа каналов.
\end{abstract}
\textbf{Ключевые слова:} теория массового обслуживания, отрицательные заявки, катастрофы, повторное обслуживание, фазовое распределение, помеченная заявка.

\begin{otherlanguage}{english}
\bigskip
\begin{center}\textbf{Calculation of Multiserver Queues with Negative Customers}\end{center}
\textbf{Abstract.} Multiserver queues with Poisson arrivals and negative events are considered. An event selects one busy server or affects the entire system. Customer removal and restart with either resampled or retained service requirements are studied. A joint phase representation is used to derive equations for the waiting and sojourn moments of a tagged customer. For selective events, the sojourn calculation accounts for changes in the number of busy servers. Retained work is represented by a finite approximation of a process with continuous coordinates, preserving dependence between servers under simultaneous restart. Phase approximation, state-space truncation, extrapolation sensitivity and simulation uncertainty are assessed separately. Numerical examples cover several loads, event rates and server counts.

\textbf{Keywords:} queueing theory, negative customers, catastrophes, service restart, phase-type distribution, tagged customer.
\end{otherlanguage}

\section{Введение}
Отрицательные заявки описывают воздействия, которые удаляют работу из системы или прерывают её выполнение. Для сетей с положительными и отрицательными заявками известны стационарные распределения в мультипликативной форме при специальных предположениях о потоках, обслуживании и механизме удаления \cite{gelenbe1991,gelenbe1992}. Обзор таких моделей приведён в \cite{artalejo2000}. Для одноканальных систем исследованы отрицательные заявки \cite{harrison1996} и мгновенное очищение при катастрофе \cite{jain1996}.

Повторное обслуживание связано с задачами первого прохождения с перезапуском \cite{pal2017}. Для системы \(M/G/1\) показано, как перезапуск меняет распределение времени завершения и характеристики очереди; при этом принципиально различаются новая случайная длительность попытки и неизменный объём задания \cite{bonomo2022}. Эти результаты для одной заявки не определяют совместное поведение нескольких каналов: общий момент прерывания связывает их состояния, а при равновероятном выборе занятого канала индивидуальная интенсивность воздействия меняется вместе с занятостью системы.

В настоящей работе эти зависимости учитываются в совместной фазовой модели. Стационарные вероятности определяются из уравнений баланса, а временные характеристики --- по уравнениям для помеченной заявки. Метод Такахаси--Таками используется как базовый подход к малым двухфазным блокам \cite{takahashi1976,ryzhikov2018ru}; расширенные модели рассчитываются разреженными линейными методами. Для сохранённой работы исследуется чувствительность конечной аппроксимации. Цель вычислительных экспериментов --- установить работоспособность расчёта в указанных режимах и отделить причины расхождений с имитацией.

\section{Постановка задачи}
Система имеет \(n\) одинаковых каналов и неограниченную очередь. Положительные заявки поступают по пуассоновскому потоку интенсивности \(\lambda>0\), отрицательные воздействия --- по независимому пуассоновскому потоку интенсивности \(\delta\ge0\). Обслуживание начинается в порядке поступления. Исходные длительности \(B\) независимы, одинаково распределены и независимы от потоков. Обозначим
\[
b_k=\mathbb E[B^k],\quad CV=\frac{\sqrt{b_2-b_1^2}}{b_1},\quad \rho=\frac{\lambda b_1}{n}.
\]
Номинальная нагрузка \(\rho\) задаёт условия эксперимента. При удалениях и повторах она не равна фактической доле занятых каналов.

Выборочное воздействие RCS равновероятно выбирает один из занятых каналов. При удалении соответствующая заявка покидает систему; общий вариант удаления, называемый катастрофой, мгновенно очищает и обслуживание, и очередь. При повторе выбранная заявка либо все обслуживаемые заявки немедленно начинают обслуживание заново, сохраняя преимущество перед ожидающими. После прерывания независимо выбирается новая длительность либо сохраняется первоначальная длительность данной заявки. В обоих случаях выполненная работа теряется. В пустой системе воздействие ничего не меняет.

Время ожидания \(W\) отсчитывается от поступления до начала обслуживания; при очищении оно заканчивается также при удалении из очереди. Время пребывания \(V\) заканчивается завершением обслуживания либо удалением. Каждая поступившая заявка учитывается один раз. Обозначим \(w_k=\mathbb E[W^k]\), \(v_k=\mathbb E[V^k]\), а через \(q\) --- вероятность завершить обслуживание. Таблица~\ref{tab:methods} связывает механизмы воздействий с расчётными моделями.

\begin{table}[!htbp]
\centering\small
\begin{tabular}{@{}p{3.6cm}p{6.0cm}p{5.0cm}@{}}
\toprule
Механизм & Переход при воздействии & Численное приближение \\
\midrule
RCS, удаление & Удаление выбранной заявки; риск для каждой равен \(\delta/M(t)\) & PH-представление \(B\), пределы \(N,H\) \\
Очищение & Переход в пустую систему & PH-представление \(B\), предел \(N\) \\
RCS, новая длительность & Сброс фазы выбранной заявки по начальному распределению & PH-представление \(B\), пределы \(N,H\) \\
Все каналы, новая длительность & Одновременный сброс всех фаз & PH-представление \(B\), предел \(N\) \\
RCS, сохранённая длительность & Сброс прогресса выбранной работы с сохранением её класса & Квадратура \(J\), фазы \(K\), пределы \(N,H\) \\
Все каналы, сохранённая длительность & Одновременный сброс прогресса всех работ с сохранением классов & Квадратура \(J\), фазы \(K\), предел \(N\) \\
\bottomrule
\end{tabular}
\caption{Расчётные модели. \(M(t)\) --- текущее число занятых каналов; \(N\) --- число стационарных уровней; \(H\) --- предел числа заявок после помеченной. При общих воздействиях отдельное усечение \(H\) не требуется при использовании формул для одной заявки.}
\label{tab:methods}
\end{table}

\section{Совместная фазовая модель}
\subsection{Представление обслуживания и переходов}
Фазовое распределение (PH) задаётся начальным вектором \(\boldsymbol\alpha\) и переходной частью генератора \(S\). Здесь \(\boldsymbol\alpha\ge0\), \(\boldsymbol\alpha\mathbf1=1\), внедиагональные элементы \(S\) неотрицательны, а \(\boldsymbol s_0=-S\mathbf1\ge0\) содержит интенсивности завершения. Преобразование Лапласа—Стилтьеса (ПЛС) и начальные моменты равны
\[
\beta(s)=\boldsymbol\alpha(sI-S)^{-1}\boldsymbol s_0,\quad
b_k=k!\boldsymbol\alpha(-S)^{-k}\mathbf1.
\]
Число заявок в системе задаёт уровень \(i\). При \(m=\min(i,n)\) дополнительно учитывается неупорядоченный набор фаз \(m\) занятых каналов. Перестановки одинаковых каналов объединяются без предположения независимости. При \(d\) фазах на насыщенном уровне имеется \(\binom{d+n-1}{n}\) состояний. Поступление в очередь сохраняет фазы; поступление на свободный канал выбирает фазу по \(\boldsymbol\alpha\). Завершение удаляет фазу, а при непустой очереди освободившийся канал сразу начинает следующее обслуживание.

Пусть \(R\) --- стохастическая матрица сброса одной фазы. Для новой длительности все её строки равны \(\boldsymbol\alpha\). До объединения перестановок добавка к генератору при \(m>0\) имеет вид
\[
G_m^{\mathrm{RCS}}=\frac{\delta}{m}\sum_{\ell=1}^{m}(R_\ell-I),\qquad
G_m^{\text{общ}}=\delta(R^{\otimes m}-I),
\]
где \(R_\ell\) действует на фазу канала \(\ell\), сохраняя остальные координаты. Общий сброс является одним событием. При удалении RCS каждый занятый канал освобождается дополнительно с интенсивностью \(\delta/m\); катастрофа переводит любое непустое состояние в пустое. Явные двухфазные блоки и вариант с вспомогательными состояниями приведены в приложении.

\subsection{Стационарные вероятности и устойчивость}
Обозначим через \(Q\) генератор совместной цепи и через \(\boldsymbol\pi\) её стационарный вектор. Уравнения расчёта имеют вид
\[
\boldsymbol\pi Q=0,\quad \boldsymbol\pi\mathbf1=1.
\]
В реализации учитываются уровни \(0,\ldots,N-1\); поступления на верхней границе подавляются. Малые системы решаются прямым разреженным методом, большие --- итерационным методом для вложенной цепи. Контролируются нормировка, неотрицательность и невязка баланса, затем отдельно увеличивается \(N\). Малая масса последнего уровня сама по себе не подтверждает точность второго момента.

Для фазовых моделей без очищения проверяется положительный запас пропускной способности насыщенной системы. Когда заняты все каналы, индивидуальная интенсивность равна \(\delta/n\) для RCS и \(\delta\) для общих прерываний. По соответствующему фазовому времени освобождения канала \(T_{\mathrm{sat}}\) проверяется \(\lambda\mathbb E[T_{\mathrm{sat}}]<n\). Это проверка выбранной конечной фазовой модели; существование моментов исходного процесса с сохранённой непрерывной работой рассматривается отдельно. При очищении с \(\delta>0\) время пребывания ограничено временем до следующей катастрофы, поэтому допускаются примеры с \(\rho>1\).

\subsection{Положительная аппроксимация при малом коэффициенте вариации}
Согласование трёх моментов двухфазным гиперэкспоненциальным выражением может давать веса вне \([0,1]\) либо комплексные параметры. Такие коэффициенты нельзя считать вероятностями и интенсивностями переходов. При этом отрицательный вес в одном представлении ещё не исключает положительного фазового представления того же преобразования. Для использованного ниже \(CV=0{,}8\) прежнее рациональное выражение эквивалентно двухфазному распределению Кокса; эта эквивалентность проверена по ПЛС и характеристикам очереди.

Для более подробной аппроксимации исходного гамма-распределения положим \(a=1/CV^2=m+r\), \(\theta=b_1/a\), где \(m\) целое и \(0<r<1\). Из интеграла Эйлера \cite{nistdlmf} следует
\[
\mathbb E\!\left[\frac{1}{1+\theta sX}\right]=(1+\theta s)^{-r},\quad
X\sim\operatorname{Beta}(r,1-r).
\]
Следовательно, \(\operatorname{Gamma}(r,\theta)\) представляется смесью экспоненциальных длительностей с условным средним \(\theta X\). Положительная квадратура Гаусса с \(J\) узлами даёт конечную смесь. Перед ней добавляются \(m\) последовательных экспоненциальных фаз среднего \(\theta\). Квадратура сохраняет первые \(2J-1\) моментов исходного гамма-распределения, поскольку условный момент порядка \(k\) пропорционален \(X^k\). При целочисленном \(a\) используется точное распределение Эрланга. Совпадение конечного числа моментов не заменяет проверку изменения характеристик очереди при увеличении \(J\).

\section{Временные характеристики}\label{sec:times}
\subsection{Уравнения для помеченной заявки}
Пуассоновское поступление наблюдает стационарное распределение совместного состояния. Пока помеченная заявка ожидает, учитываются число заявок перед ней, число заявок после неё и фазы занятых каналов. После начала обслуживания учитываются её фаза, число остальных заявок и фазы остальных занятых каналов. Для времени ожидания поглощением служит начало обслуживания или удаление из очереди; для времени пребывания --- завершение обслуживания или удаление помеченной заявки.

Пусть \(\mathcal T\) --- переходная часть генератора такого процесса. Векторы условных начальных моментов времени до поглощения удовлетворяют уравнениям первого шага
\[
(-\mathcal T)\boldsymbol u_k=k\boldsymbol u_{k-1},\quad
\boldsymbol u_0=\mathbf1,\quad k\ge1.
\]
Усреднение по состоянию, наблюдаемому приходом, даёт \(w_k\) или \(v_k\). Уравнения сохраняют изменение индивидуального риска \(\delta/M(t)\) при RCS. В частности, для удаления не требуется фиксировать число занятых каналов на момент начала обслуживания. Для повторов не вводится усреднённая интенсивность, зависящая от эффективных моментов.

Поступления после помеченной заявки не влияют на её начальное ожидание, но могут изменить риск RCS во время её обслуживания. Поэтому при расчёте \(V\) они учитываются до отдельного предела \(H\), на котором дальнейшие поступления подавляются. Усечения \(N\) и \(H\) проверяются раздельно. Моменты фазовых времён получаются из линейных уравнений, без численного дифференцирования ПЛС. Для исходной сохранённой работы при общих прерываниях ниже используются также аналитические выражения.

Независимый контроль средних даёт закон Литтла, причём учитываются все заявки, включая удалённые:
\[
\lambda w_1=\sum_i\max(i-n,0)p_i,\qquad
\lambda v_1=\sum_i i p_i,
\]
где \(p_i\) --- суммарная вероятность уровня. При удалении RCS баланс потока даёт \(q=1-\delta(1-p_0)/\lambda\). При очищении \(q\) определяется отношением интенсивности успешных завершений к \(\lambda\); без удалений в устойчивом режиме \(q=1\).

\subsection{Сохранённая работа}
При повторе первоначальной длительности точное состояние содержит число ожидающих и пары \((B_\ell,a_\ell)\) для занятых каналов. Здесь \(B_\ell\) --- исходный объём работы, \(a_\ell\) --- выполненная часть. Между событиями \(da_\ell/dt=1\), при \(a_\ell=B_\ell\) обслуживание завершается. RCS с интенсивностью \(\delta/M(t)\) заменяет выбранную координату \(a_\ell\) нулём; общее прерывание одновременно обнуляет все координаты прогресса. Значения \(B_\ell\) сохраняются. Это марковский процесс с непрерывными координатами и детерминированным движением между скачками.

Для численного решения закон \(B\) заменяется положительной квадратурой с узлами \(d_h\) и весами \(\omega_h\), \(h=1,\ldots,J\). Класс работы выбирается один раз. Его выполнение приближается \(K\) последовательными экспоненциальными фазами интенсивности \(K/d_h\). Матрица \(R\) возвращает заявку в первую фазу того же класса. При конечном \(K\) длительность прохождения фаз случайна и имеет дисперсию \(d_h^2/K\); при \(K\to\infty\) она сходится к \(d_h\). Таким образом, конечная модель приближает выполнение сохранённой работы, а не воспроизводит его точно.

Совместное состояние содержит неупорядоченный набор пар «класс, фаза». Для RCS моменты \(W,V\) находятся по полным уравнениям помеченной заявки. При общих прерываниях начальное ожидание \(W\) независимо от будущего времени завершения собственной работы \(T\): её объём независим от предшествующего ожидания, а будущие приращения пуассоновского потока независимы от прошлого. Поэтому
\[
v_1=w_1+\mathbb E[T],\qquad
v_2=w_2+2w_1\mathbb E[T]+\mathbb E[T^2].
\]
Это не означает независимости разных каналов. Для RCS независимость \(W\) и \(T\) не предполагается.

При общих прерываниях моменты \(T\) для исходного закона определяются через \(M_B(t)=\mathbb E[e^{tB}]\):
\[
\mathbb E[T]=\frac{M_B(\delta)-1}{\delta},\qquad
\mathbb E[T^2]=\frac{2}{\delta^2}\bigl(M_B(2\delta)-M_B(\delta)-\delta M'_B(\delta)\bigr).
\]
Эти выражения используются только в области существования моментов; при \(\delta=0\) берутся предельные значения. Для гамма-распределения со шкалой \(\theta\) необходимы соответственно \(\delta\theta<1\) и \(2\delta\theta<1\), для положительного \(H_2\) --- \(\delta<\min\mu_\ell\) и \(2\delta<\min\mu_\ell\) по фазам положительного веса. В экспериментах с сохранённой работой дополнительно проверяется более сильное условие конечности пятого момента времени завершения при постоянной интенсивности \(\delta\). Оно используется как консервативная проверка хвоста времени завершения одной заявки: интенсивность RCS не превосходит \(\delta\). Это не общий критерий существования стационарных моментов очереди.

Число классов \(J\), число фаз \(K\), уровни \(N\) и предел \(H\) уточняются отдельно. Экстраполяция применяется к начальным моментам; стандартные отклонения вычисляются после неё. Коэффициенты экстраполяции определяются выбранными сетками, без использования имитационных значений. Сходимость времени прохождения одной работы сама по себе не доказывает сходимость стационарных моментов очереди; разности сеток ниже рассматриваются как численная диагностика.
